%\section{PAGINA 8}

\begin{table}[h!]
    \centering
    \caption{Resultados del experimento.}
    \label{tab:27-1}
    \begin{tabular}{|c|c|c|c|c|c|c|c|}
        \hline
        \multirow{3}{*}{$R_L$} & \multicolumn{2}{c|}{$I_N$ (\si{\milli\ampere})} & \multicolumn{2}{c|}{$R_N$ (\si{\ohm})} & \multicolumn{3}{c|}{$I_L$ (\si{\milli\ampere})} \\
        \cline{2-8}
        & \multirow{2}{*}{Medida} & \multirow{2}{*}{Calculada} & \multirow{2}{*}{Medida} & \multirow{2}{*}{Calculada} & \multicolumn{2}{c|}{Medida} & \multirow{2}{*}{Calculada} \\
        \cline{6-7}
        & & & & & Norton & Circuito original & \\
        \hline
        \SI{5100}{\ohm} & & & & & & & \\
        \hline
         & & & & & & & \\
        \hline
         & & & & & & & \\
        \hline
         & & & & & & & \\
        \hline
    \end{tabular}
\end{table}

\vspace{2em}

\begin{enumerate}[label=\arabic*., resume]
    \setcounter{enumi}{6} % Continuamos la lista en el número 7
    \item Si se utilizan en el experimento valores de $R_3$ (fig. 27-6 a) distintos de \SI{5100}{\ohm}, ¿cuál debe ser el valor de $I_N$ en la figura 27-6 b? Explicarlo.
\end{enumerate}

\vspace{1cm}

\section{Respuestas a las preguntas de autoexamen}
\begin{enumerate}[label=\arabic*.]
    \item 
    \begin{enumerate}[label=\alph*)]
        \item 0,3
        \item 240
        \item 0,141
    \end{enumerate}
    \item 
    \begin{enumerate}[label=\alph*)]
        \item 0,05
        \item 346
        \item 0,021
    \end{enumerate}
\end{enumerate}